%!TEX root = ../thesis.tex
%*******************************************************************************
%****************************** Chương 2 ***************************************
%*******************************************************************************

\chapter{CÁC NGHIÊN CỨU LIÊN QUAN}
\label{chap:related}

\ifpdf
    \graphicspath{{Chapter2/Figs/Raster/}{Chapter2/Figs/PDF/}{Chapter2/Figs/}}
\else
    \graphicspath{{Chapter2/Figs/Vector/}{Chapter2/Figs/}}
\fi

\nomenclature[z-OCR-VQA]{Scene-text VQA}{VQA mà câu trả lời phải đọc từ chữ xuất hiện trong ảnh}
\nomenclature[z-Q-Former]{Q-Former}{Querying Transformer: bộ nối ảnh--ngôn ngữ của BLIP-2}

Chương này tổng quan năm nhóm công trình mà khóa luận dựa vào: bài toán VQA và
VQA đọc chữ trong ảnh (mục~\ref{sec:rel-vqa}), các vision-language model
(mục~\ref{sec:rel-vlm}), các phương pháp parameter-efficient fine-tuning
(mục~\ref{sec:rel-peft}), các nghiên cứu VQA tiếng Việt
(mục~\ref{sec:rel-vivqa}) và cách đo hallucination (mục~\ref{sec:rel-halluc}).
Mục~\ref{sec:rel-summary} tổng hợp các khoảng trống G1--G4 đã nêu ở
Chương~\ref{chap:intro} và chỉ ra khóa luận giải quyết khoảng trống nào.

%********************************** Mục 2.1 ************************************
\section{Hỏi đáp trên ảnh và VQA đọc chữ trong ảnh}
\label{sec:rel-vqa}

\subsection{Bài toán VQA}

Antol và cộng sự~\cite{antol2015vqa} đặt bài toán VQA dưới dạng mở: cho một ảnh và một câu hỏi
tự nhiên, hệ thống trả lời bằng một từ hoặc cụm từ ngắn. Bộ dữ liệu VQA của họ
thu mười câu trả lời của người cho mỗi câu hỏi và định nghĩa thước đo VQA
Accuracy: một câu trả lời được tính đúng hoàn toàn nếu ít nhất ba người cũng trả
lời như vậy, tức $\mathrm{Acc}(a) = \min(\#\{\text{người trả lời } a\}/3,\ 1)$.
Các hệ thống thời kỳ này ghép một mạng CNN mã hóa ảnh với một mạng LSTM mã hóa
câu hỏi, rồi \emph{phân loại} câu trả lời trên một tập đáp án thường gặp cố định.
Cách đặt bài toán phân loại này hiệu quả với câu trả lời ngắn và lặp lại (màu
sắc, số lượng, có/không), nhưng không thể sinh một câu trả lời chưa từng có trong
tập đáp án lúc huấn luyện.

\subsection{VQA đọc chữ trong ảnh}

Nhiều câu hỏi thực tế chỉ trả lời được bằng cách đọc chữ trong ảnh: tên cửa hàng,
giá tiền, số điện thoại. Singh và cộng sự~\cite{singh2019textvqa} giới thiệu bộ TextVQA cho
loại câu hỏi này và cho thấy các mô hình VQA thông thường gần như không trả lời
được; mô hình LoRRA của họ đưa thêm các từ OCR vào mô hình và cho phép
\emph{chép} một từ OCR làm câu trả lời. Biten và cộng sự~\cite{biten2019stvqa} xây bộ ST-VQA với
cùng mục tiêu và đề xuất thước đo ANLS (Average Normalized Levenshtein
Similarity), tính độ giống theo ký tự giữa câu trả lời và đáp án để không phạt
nặng những lỗi đọc nhỏ do OCR. M4C~\cite{hu2020m4c} mở rộng ý tưởng chép bằng một
transformer đa phương thức sinh câu trả lời từng bước, mỗi bước chọn giữa một từ
trong từ điển và một từ OCR.

Bài học chung của nhóm công trình này là: với VQA đọc chữ, chất lượng câu trả lời
bị chặn bởi chất lượng chữ đọc được. Nếu OCR đọc sai hoặc sắp sai thứ tự các vùng
chữ, mô hình không thể chép ra đáp án đúng. Đây là lý do khóa luận đo riêng chất
lượng OCR tiếng Việt và thứ tự đọc (Chương~\ref{chap:methods}).

%********************************** Mục 2.2 ************************************
\section{Vision-language model}
\label{sec:rel-vlm}

Vision-language model (VLM) hiện nay có chung một khuôn: một bộ mã hóa ảnh biến
ảnh thành chuỗi vector, một bộ nối chuyển các vector đó sang không gian của mô
hình ngôn ngữ, và một mô hình ngôn ngữ lớn (LLM) đọc chuỗi token ảnh cùng câu hỏi
để sinh câu trả lời. Khác biệt giữa các VLM nằm ở bộ nối, ở cách xử lý độ phân
giải ảnh và ở dữ liệu huấn luyện. Bảng~\ref{tab:rel-vlm} so sánh bốn VLM mà đề
cương cho phép chọn làm baseline.

\paragraph{BLIP-2.} Li và cộng sự~\cite{li2023blip2} giữ nguyên (đóng băng) cả bộ mã hóa ảnh
lẫn LLM và chỉ huấn luyện một mô-đun nhỏ ở giữa là Q-Former. Q-Former dùng một
số ít query học được để rút thông tin từ ảnh thành một chuỗi token cố định đưa
vào LLM. Cách làm này rẻ về huấn luyện, nhưng số token ảnh cố định và nhỏ khiến
mô hình khó giữ chi tiết nhỏ như chữ trong ảnh.

\paragraph{LLaVA và LLaVA-1.5.} Liu và cộng sự~\cite{liu2023llava} nối bộ mã hóa ảnh với LLM
chỉ bằng một lớp chiếu và huấn luyện trên dữ liệu hội thoại về ảnh (visual
instruction tuning). LLaVA-1.5~\cite{liu2024llava15} thay lớp chiếu bằng một MLP,
tăng độ phân giải ảnh và thêm một câu nhắc yêu cầu trả lời ngắn; riêng câu nhắc
này đã đủ để mô hình trả lời đúng kiểu ngắn mà các bộ VQA cần. LLaVA-1.5 cũng
công bố cả bản fine-tune toàn bộ và bản LoRA (mục~\ref{sec:rel-peft}).

\paragraph{Florence-2.} Xiao và cộng sự~\cite{xiao2024florence2} xây một mô hình thị giác hợp
nhất kiểu encoder--decoder, nhận yêu cầu bằng văn bản và sinh kết quả cho các tác
vụ chú thích ảnh, phát hiện vật thể, định vị và OCR. Mô hình rất nhỏ (0,23 và 0,77
tỉ tham số) nên dễ huấn luyện, nhưng không phải mô hình hội thoại đa ngôn ngữ và
không được huấn luyện cho hỏi đáp mở bằng tiếng Việt.

\paragraph{Qwen2.5-VL.} Bai và cộng sự~\cite{bai2025qwen25vl} dùng một bộ mã hóa ViT nhận ảnh ở
độ phân giải gốc (native dynamic resolution): ảnh được cắt thành các ô 14×14 điểm
ảnh, rồi bộ nối gộp mỗi nhóm 2×2 ô thành một token, nên số token ảnh tỉ lệ với
kích thước ảnh. Mô hình ngôn ngữ là Qwen2.5, đa ngôn ngữ, và mô hình được huấn
luyện nhiều trên dữ liệu tài liệu và chữ trong ảnh; bản 3 tỉ tham số đạt 79,3
trên TextVQA theo báo cáo của tác giả. Khóa luận chọn Qwen2.5-VL-3B-Instruct làm
baseline; kiến trúc chi tiết và lý do chọn được trình bày ở
Chương~\ref{chap:methods}.

\begin{table}[htbp]
\caption{So sánh bốn VLM trong phạm vi đề cương}
\label{tab:rel-vlm}
\centering
\small
\begin{tabular}{p{0.13\textwidth}p{0.25\textwidth}p{0.24\textwidth}p{0.26\textwidth}}
\toprule
\textbf{Mô hình} & \textbf{Bộ nối ảnh--ngôn ngữ} & \textbf{Điểm mạnh với đề tài} & \textbf{Hạn chế với đề tài} \\
\midrule
BLIP-2 & Q-Former, số token ảnh cố định; ảnh và LLM đóng băng & Rẻ khi huấn luyện; là nền của ViTextBLIP-2 & Ít token ảnh, dễ mất chi tiết chữ nhỏ \\
LLaVA-1.5 & MLP chiếu từng patch sang LLM & Kiến trúc đơn giản; có bằng chứng LoRA ngang fine-tune toàn bộ & Chủ yếu tiếng Anh; độ phân giải ảnh cố định \\
Florence-2 & Encoder--decoder hợp nhất, yêu cầu bằng văn bản & Rất nhỏ; có tác vụ OCR riêng & Không phải mô hình hội thoại đa ngôn ngữ \\
Qwen2.5-VL & ViT độ phân giải gốc, gộp 2×2 ô thành một token & Đa ngôn ngữ; đọc chữ trong ảnh tốt; có bản 3B & Số token ảnh tăng theo kích thước ảnh, tốn bộ nhớ \\
\bottomrule
\end{tabular}
\end{table}

%********************************** Mục 2.3 ************************************
\section{Parameter-efficient fine-tuning}
\label{sec:rel-peft}

\subsection{Vì sao cần PEFT}

Fine-tune toàn bộ một mô hình lớn cho mỗi bài toán nghĩa là cập nhật và lưu một
bản sao của mọi tham số. Với GPT-3 175 tỉ tham số, mỗi bản sao chiếm khoảng 350~GB
ở 16-bit, và lúc huấn luyện còn phải giữ thêm gradient và trạng thái của bộ tối
ưu~\cite{hu2022lora}. Parameter-efficient fine-tuning (PEFT) giữ nguyên trọng số
gốc và chỉ huấn luyện một phần rất nhỏ tham số thêm vào. Các hướng sớm gồm
adapter~\cite{houlsby2019adapter}, chèn các lớp nhỏ vào giữa các lớp của mạng, và
prefix-tuning~\cite{li2021prefix}, học một đoạn vector đặt trước chuỗi đầu vào.
Adapter làm tăng độ trễ khi suy luận vì thêm lớp tuần tự; prefix-tuning chiếm
một phần độ dài chuỗi và khó tối ưu.

\subsection{LoRA}

Hu và cộng sự~\cite{hu2022lora} giả định rằng phần thay đổi của trọng số khi thích nghi có hạng
thấp. Với một ma trận trọng số $W_0 \in \mathbb{R}^{d \times k}$, LoRA giữ $W_0$
cố định và học hai ma trận nhỏ $B \in \mathbb{R}^{d \times r}$,
$A \in \mathbb{R}^{r \times k}$ với $r \ll \min(d, k)$:
\begin{equation}
  h = W_0 x + \frac{\alpha}{r} B A x .
  \label{eq:lora}
\end{equation}
$B$ được khởi tạo bằng 0 nên lúc bắt đầu mô hình giữ nguyên hành vi gốc. Sau khi
huấn luyện, $BA$ có thể cộng thẳng vào $W_0$, nên suy luận không chậm hơn. Trên
GPT-3, LoRA đạt kết quả ngang hoặc hơn fine-tune toàn bộ với số tham số huấn
luyện ít hơn khoảng 10.000 lần và bộ nhớ GPU giảm khoảng 3 lần.

\subsection{QLoRA}

LoRA giảm số tham số phải học nhưng vẫn phải giữ trọng số gốc ở 16-bit trong bộ
nhớ. Dettmers và cộng sự~\cite{dettmers2023qlora} lượng tử hóa trọng số gốc xuống 4-bit và huấn
luyện LoRA ở trên, với ba kỹ thuật: kiểu số 4-bit NormalFloat (NF4) được thiết
kế cho trọng số có phân phối gần chuẩn; lượng tử hóa kép (double quantization)
nén cả các hằng số lượng tử hóa; và bộ tối ưu phân trang (paged optimizer) để
tránh tràn bộ nhớ khi có đỉnh tiêu thụ. Nhờ đó, một LLM 65 tỉ tham số được
fine-tune trên một GPU 48~GB, trong khi kết quả khớp với fine-tune 16-bit trên các
thí nghiệm của tác giả. Bài báo cũng khuyến nghị gắn LoRA lên mọi lớp tuyến tính
để bù sai số lượng tử hóa. Khóa luận dùng QLoRA vì đây là cách duy nhất để
fine-tune một VLM vài tỉ tham số trên một GPU Colab; lựa chọn gắn LoRA chỉ vào các
ma trận attention hay vào mọi lớp tuyến tính được kiểm tra bằng ablation ở
Chương~\ref{chap:results}.

\subsection{PEFT cho vision-language model}

Sung và cộng sự~\cite{sung2022vladapter} so sánh có kiểm soát các phương pháp PEFT trên các bài
toán ảnh--ngôn ngữ, trong đó có VQA, và cho thấy một adapter dùng chung chỉ với
4,18\% tham số huấn luyện đạt kết quả gần bằng fine-tune toàn bộ. LLaVA-1.5 cho
bằng chứng tương tự ở quy mô lớn hơn: với mô hình 13 tỉ tham số, bản LoRA đạt
80,0 trên VQAv2 (bằng bản fine-tune toàn bộ), 60,2 so với 61,3 trên TextVQA, và
86,7 so với 85,9 F1 trên POPE~\cite{liu2024llava15}. Bảng~\ref{tab:rel-peft} tóm
tắt các cách thích nghi.

\begin{table}[htbp]
\caption{Các cách thích nghi một mô hình lớn}
\label{tab:rel-peft}
\centering
\small
\begin{tabular}{p{0.20\textwidth}p{0.22\textwidth}p{0.20\textwidth}p{0.26\textwidth}}
\toprule
\textbf{Cách} & \textbf{Tham số huấn luyện} & \textbf{Độ trễ suy luận} & \textbf{Bộ nhớ khi huấn luyện} \\
\midrule
Fine-tune toàn bộ & Toàn bộ & Không đổi & Rất lớn: trọng số, gradient, trạng thái tối ưu \\
Adapter & Vài phần trăm & Tăng (thêm lớp) & Giữ trọng số 16-bit \\
Prefix-tuning & Rất ít & Tăng nhẹ (chuỗi dài hơn) & Giữ trọng số 16-bit \\
LoRA & Dưới 1\% & Không đổi sau khi gộp & Giữ trọng số 16-bit \\
QLoRA & Dưới 1\% & Không đổi sau khi gộp & Trọng số 4-bit; nhỏ nhất \\
\bottomrule
\end{tabular}
\end{table}

Tuy vậy, các bằng chứng trên đều đến từ tiếng Anh. Trong các công trình nhóm khảo
sát, chưa có nghiên cứu nào kiểm chứng LoRA hoặc QLoRA trên VLM cho VQA tiếng Việt
hay cho VQA đọc chữ tiếng Việt. Đây là khoảng trống G3 mà khóa luận tập trung.

%********************************** Mục 2.4 ************************************
\section{VQA tiếng Việt}
\label{sec:rel-vivqa}

\subsection{Bộ dữ liệu}

\paragraph{ViVQA.} Tran và cộng sự~\cite{tran2021vivqa} công bố bộ dữ liệu VQA tiếng Việt đầu
tiên, gồm 10.328 ảnh và 15.000 cặp câu hỏi--câu trả lời. Câu trả lời ngắn, phần
lớn chỉ một đến ba từ, nên bài toán thường được đặt dưới dạng phân loại.

\paragraph{OpenViVQA.} Nguyen và cộng sự~\cite{nguyen2023openvivqa} xây bộ dữ liệu với câu trả lời
mở, dài hơn và đòi hỏi sinh câu trả lời thay vì chọn từ một tập cố định; khoảng
một nửa số mẫu có liên quan tới chữ trong ảnh~\cite{nguyen2024vitextvqa}.

\paragraph{ViTextVQA.} Nguyen và cộng sự~\cite{nguyen2024vitextvqa} công bố bộ dữ liệu VQA đọc chữ
tiếng Việt quy mô lớn đầu tiên, gồm 16.762 ảnh và 50.342 câu hỏi, chia train, dev,
test theo tỉ lệ 7:1:2. Câu trả lời thường là chữ chép từ ảnh như tên cửa hàng,
khẩu hiệu, giá tiền. Người đạt 88,00 EM và 96,42 F1 trên bộ này. Các tác giả cũng
cho thấy sắp xếp chữ OCR theo thứ tự từ trái sang phải, từ trên xuống dưới giúp
các mô hình trả lời tốt hơn. Khóa luận dùng ViVQA và ViTextVQA; thống kê chi tiết
được trình bày ở Chương~\ref{chap:methods}.

\subsection{Phương pháp}

Bảng~\ref{tab:rel-vivqa} tóm tắt các phương pháp tiêu biểu theo thời gian.

\begin{table}[htbp]
\caption{Các nghiên cứu VQA tiếng Việt tiêu biểu}
\label{tab:rel-vivqa}
\centering
\small
\begin{tabular}{p{0.21\textwidth}p{0.33\textwidth}p{0.16\textwidth}p{0.20\textwidth}}
\toprule
\textbf{Nghiên cứu} & \textbf{Cách làm} & \textbf{Kết quả công bố} & \textbf{Hạn chế} \\
\midrule
ViVQA~\cite{tran2021vivqa}, 2021 & LSTM, Bi-LSTM với nhiều kiểu word embedding; Hierarchical Co-Attention; chọn đáp án trong tập cố định & Acc 34,96 trên ViVQA & Đặc trưng rời rạc; không đọc được chữ \\
Bi-directional Cross-Attention~\cite{nguyentran2022bica}, 2022 & Bộ mã hóa đã tiền huấn luyện, cross-attention hai chiều giữa ảnh và câu hỏi & Vượt baseline ViVQA & Vẫn là bài toán phân loại \\
PAT~\cite{nguyen2023pat}, 2023 & Parallel Attention Transformer, trích đặc trưng câu hỏi theo tầng & Acc 60,55 trên ViVQA & Phân loại; không đọc chữ \\
BARTPhoBEiT~\cite{tran2023bartphobeit}, 2023 & BARTpho cho tiếng Việt kết hợp bộ mã hóa ảnh kiểu BEiT & Acc 68,58 trên ViVQA & Phân loại; không đọc chữ \\
AViVQA-TranConI~\cite{nguyen2024tranconi}, 2024 & BLIP-2 (đặc trưng toàn cục) và EfficientNet (đặc trưng cục bộ), BARTpho, trộn bằng BEiT-3; phân loại trên 353 đáp án & Acc 71,04 trên ViVQA & Không sinh được đáp án mới; không có OCR \\
ViTextBLIP-2~\cite{nguyen2024vitextvqa}, 2025 & SwinTextSpotter đọc chữ, Q-Former nối ảnh và chữ OCR, ViT5 đóng băng sinh câu trả lời & EM 25,48; F1 53,95 trên ViTextVQA & Phụ thuộc OCR ngoài; cách xa người (88,00 EM) \\
\bottomrule
\end{tabular}
\end{table}

Các nghiên cứu trên đi theo một lộ trình rõ: từ đặc trưng rời rạc kết hợp
attention (2021), sang bộ mã hóa đã tiền huấn luyện và mô-đun trộn đa phương thức
(2022--2024), rồi tới mô hình sinh dựa trên BLIP-2 cho dữ liệu có chữ (2025). Có
hai vướng mắc lặp lại. Thứ nhất, phần lớn các công trình trên ViVQA coi VQA là
phân loại trên tập đáp án cố định, nên không trả lời được những đáp án chưa thấy
khi huấn luyện, điều bắt buộc với câu hỏi đọc chữ. Thứ hai, đọc chữ tiếng Việt có
dấu trong ảnh vẫn là điểm yếu: ViTextBLIP-2 dựa vào một bộ OCR bên ngoài được giữ
cố định khi huấn luyện, nên lỗi đọc chữ của OCR truyền thẳng vào câu trả lời.

Cần lưu ý khi so sánh: Accuracy của các phương pháp phân loại trên ViVQA được tính
trên tập đáp án cố định, không cùng cách chấm với EM của mô hình sinh tự do mà
khóa luận dùng. Vì vậy các con số trên ViVQA ở Bảng~\ref{tab:rel-vivqa} chỉ để
tham khảo. Trên ViTextVQA, khóa luận so trực tiếp với ViTextBLIP-2 vì cùng tập
test và cùng thước đo EM, F1 (Chương~\ref{chap:results}).

%********************************** Mục 2.5 ************************************
\section{Hallucination và cách đo}
\label{sec:rel-halluc}

Hallucination trong VLM là hiện tượng mô hình sinh nội dung trôi chảy nhưng không
có trong ảnh, ví dụ nhắc tới một vật không xuất hiện. Rohrbach và cộng sự~\cite{rohrbach2018chair}
đo hiện tượng này trên chú thích ảnh bằng thước đo CHAIR, đếm tỉ lệ vật được nhắc
tới mà không có trong nhãn của ảnh. Cách đo này phụ thuộc vào việc tách tên vật ra
khỏi câu văn và vào độ dài câu.

POPE~\cite{li2023pope} đổi sang dạng câu hỏi có/không: với mỗi ảnh, hỏi ``Trong
ảnh có \emph{vật X} không?'' cho cả vật có thật và vật không có. Vật không có được
chọn theo ba cách: ngẫu nhiên, vật phổ biến trong toàn bộ dữ liệu, và vật hay xuất
hiện cùng các vật trong ảnh (adversarial). Kết quả báo cáo bằng Accuracy,
Precision, Recall, F1 và tỉ lệ trả lời ``có''. Vì câu trả lời chỉ là có hoặc
không, POPE không cần phân tích câu văn và dễ áp dụng cho ngôn ngữ khác. Tuy vậy,
chưa có bộ đo kiểu POPE nào bằng tiếng Việt (khoảng trống G4). Khóa luận xây bộ
POPE-vi theo đúng giao thức này để kiểm tra giả thuyết H4
(Chương~\ref{chap:discussion}).

%********************************** Mục 2.6 ************************************
\section{Tổng kết và khoảng trống nghiên cứu}
\label{sec:rel-summary}

Bảng~\ref{tab:rel-gaps} nối mỗi khoảng trống nêu ở Chương~\ref{chap:intro} với
bằng chứng từ các công trình trên và với cách khóa luận xử lý.

\begin{table}[htbp]
\caption{Khoảng trống, bằng chứng và cách khóa luận xử lý}
\label{tab:rel-gaps}
\centering
\small
\begin{tabular}{p{0.20\textwidth}p{0.36\textwidth}p{0.34\textwidth}}
\toprule
\textbf{Khoảng trống} & \textbf{Bằng chứng} & \textbf{Khóa luận xử lý} \\
\midrule
G1. Tiền huấn luyện thiên tiếng Anh & Các VLM chủ yếu huấn luyện trên tiếng Anh; LLaVA-1.5 chỉ có khả năng đa ngôn ngữ nổi lên, không được huấn luyện riêng & Chọn baseline đa ngôn ngữ; đo lỗi sai ngôn ngữ, sai định dạng ở chế độ zero-shot \\
G2. VQA tiếng Việt yếu, nhất là đọc chữ & ViTextBLIP-2 đạt 25,48 EM so với 88,00 của người; các phương pháp trên ViVQA là phân loại, không đọc chữ & Sinh câu trả lời tự do; đưa chữ OCR có dấu, đúng thứ tự vào prompt \\
G3. Chưa có PEFT cho VLM tiếng Việt & LoRA, QLoRA, VL-Adapter, LLaVA-1.5 chỉ kiểm chứng trên tiếng Anh & QLoRA trên Qwen2.5-VL-3B; QLoRA nhận biết OCR là đóng góp chính \\
G4. Chưa có bộ đo hallucination tiếng Việt & POPE và CHAIR chỉ có cho tiếng Anh & Xây POPE-vi; kiểm tra QLoRA có làm tăng hallucination không \\
\bottomrule
\end{tabular}
\end{table}

Tóm lại, các công trình trước cho thấy VLM đọc được chữ trong ảnh, PEFT đủ để
thích nghi VLM, và thứ tự chữ OCR quan trọng với VQA tiếng Việt. Điều còn thiếu là
một nghiên cứu ghép ba ý này cho tiếng Việt dưới ngân sách tài nguyên nhỏ, có
baseline tái lập, ablation và đo hallucination. Chương~\ref{chap:methods} trình
bày cách khóa luận làm điều đó.
